\section{Discussion}
% =============================================================================


\paragraph{A unified framework.}
The committed-action and clock mechanisms are two realizations of the
same principle: planning should be cost-aware by construction. In
committed-action environments, cost-awareness is built into the MCTS
tree, which searches over the future state induced by planning delay. In
clock environments, it is built into the value function, which is
trained on clock-augmented observations.
The gating policy can exploit both signals using the same learning algorithm.

\paragraph{Deployment as a test of training fidelity.}
The two-GPU deployment result is not just an engineering demonstration.
It shows that the committed-action abstraction, in which $K{-}1$ reflex
actions fill the planning delay during training, is a good model of
true asynchronous execution. An agent trained under this abstraction
transfers directly to hardware-separated deployment.

\paragraph{Limitations.}
The gating policy is trained on top of a fixed base planner, so joint
optimization remains open. We also study a discrete budget set
$K \in \Kset$; richer compute vocabularies, and eventually continuous
allocation, are natural extensions. Finally, our environments are
either simulated or played at controlled speeds; physical real-time
settings with uncontrolled latency would require adaptive calibration of
the timing model.


% =============================================================================
\section{Conclusion}
% =============================================================================

We studied planning in real-time reinforcement learning, where the cost
of deliberation is paid through game outcomes rather than wall-clock
seconds alone. A lightweight gating policy trained on top of a frozen
planner learns to allocate computation by state, spending more on
constrained or high-stakes decisions and less on easy ones. Across five
environments spanning two real-time mechanisms, this policy
consistently beats fixed-budget baselines while often using less
average compute. Adaptive, state-conditioned compute allocation is a
practical way to make planning-based RL systems more capable and more
efficient in real-time settings.


\begin{ack}
\todo{Acknowledgments, funding disclosure, and competing interests.}
\end{ack}


% =============================================================================
% =============================================================================

\bibliographystyle{plainnat}
\bibliography{refs}


%%%%%%%%%%%%%%%%%%%%%%%%%%%%%%%%%%%%%%%%%%%%%%%%%%%%%%%%%%%%
